\documentclass[12pt]{article}
\usepackage[T1]{fontenc}
\usepackage[utf8]{inputenc}
\usepackage{lmodern}
\usepackage{textcomp}
\usepackage{enumitem}
\usepackage{geometry}
\geometry{margin=1in}
\begin{document}
\begin{center}
Answer Key - Sociology - Graduate Level Assessment 3 \
\small (Answers are ordered exactly as in the question paper.)
\end{center}

\section*{Multiple Choice}
\begin{enumerate}[label=\arabic*.]
\item \textbf{Answer:} B
\\ \textit{Options:} A) charismatic authority; B) rational-legal authority; C) traditional authority; D) value-rational authority
\\ \textit{Note:} Weber argued that rational-legal authority legitimizes the state's monopoly on the use of force because it is based on established laws and procedures, which create a framework of legitimacy and acceptance among the populace.
\item \textbf{Answer:} A
\\ \textit{Options:} A) human behaviour is meaningful, and varies between individuals and cultures; B) it is difficult for sociologists to gain access to a research laboratory; C) sociologists are not rational or critical enough in their approach; D) we cannot collect empirical data about social life
\\ \textit{Note:} The correct answer highlights that human behavior is influenced by complex social, cultural, and individual meanings, making it fundamentally different from the objective and predictable phenomena studied in the natural sciences.
\item \textbf{Answer:} C
\\ \textit{Options:} A) secondary socialization and strict discipline; B) emotional support and sexual gratification; C) primary socialization and personality stabilization; D) oppressing women and reproducing the labour force
\\ \textit{Note:} Parsons identified primary socialization as the process through which children learn the norms and values of their society, and personality stabilization as the family's role in providing emotional support and stability for adults, making these two functions essential for the functioning of modern society.
\item \textbf{Answer:} A
\\ \textit{Options:} A) If all of your friends jumped off a bridge, I suppose you would too.; B) You should be proud to be a part of this organization.; C) If the door is closed, try the window.; D) Once a thief, always a thief.
\\ \textit{Note:} This answer illustrates the key concept of differential association theory, which posits that individuals learn behaviors and values through their associations with others, suggesting that peer influence is a significant factor in shaping one's actions and choices.
\item \textbf{Answer:} B
\\ \textit{Options:} A) indoctrination and propaganda; B) freedom of movement for citizens; C) one-party rule; D) a centrally planned economy
\\ \textit{Note:} Totalitarian societies are characterized by centralized control and repression of individual freedoms, including restrictions on citizens' movement, making "freedom of movement for citizens" an incorrect characteristic of such societies.
\end{enumerate}

\section*{True / False}
\begin{enumerate}[label=\arabic*.]
\setcounter{enumi}{5}
\item \textbf{Answer:} False
\\ \textit{Options:} A) True; B) False
\\ \textit{Note:} Stone's research indicates that variations of family structures, including nuclear families, have existed in various forms throughout history, predating industrialization.
\item \textbf{Answer:} False
\\ \textit{Options:} A) True; B) False
\\ \textit{Note:} Sullivan's study found that men tended to perform less housework when their wives were present, as the expectation for shared household responsibilities was often undermined by traditional gender roles.
\item \textbf{Answer:} True
\\ \textit{Options:} A) True; B) False
\\ \textit{Note:} The statement is true because the heightened fear and societal response to HIV and AIDS emerged in the 1980 s as a reaction to the sexual revolution of the 1960 s, which had previously challenged traditional sexual norms and behaviors.
\item \textbf{Answer:} True
\\ \textit{Options:} A) True; B) False
\\ \textit{Note:} Domhoff's analysis of power structures in the USA emphasizes the role of elite groups in maintaining control and influence, rather than framing exploitation as a decision-making process.
\item \textbf{Answer:} False
\\ \textit{Options:} A) True; B) False
\\ \textit{Note:} Smart's work highlights that while women may engage in violent crime, the overall rates and patterns of violent crime are significantly lower for women compared to men, thereby reinforcing traditional gender stereotypes about aggression and criminal behavior.
\end{enumerate}

\section*{Long Form Response}
\begin{enumerate}[label=\arabic*.]
\setcounter{enumi}{10}
\item \textbf{Answer:} The functionalist perspective in sociology posits that society is a complex system whose parts work together to promote stability and order. It explains that various social institutions, such as family, education, and religion, serve essential functions that contribute to societal cohesion and the maintenance of norms and values. One strength of functionalism lies in its ability to highlight the interconnectedness of social structures and their roles in promoting social stability. However, its limitations include a tendency to overlook social change and conflict, as well as an insufficient focus on power dynamics and inequalities that can disrupt social harmony. Overall, while functionalism provides valuable insights into the mechanisms of social order, it may fall short in addressing the complexities and challenges of contemporary society.
\\ \textit{Note:} The answer correctly identifies functionalism as a perspective that views society as a system of interconnected parts working together to maintain order and stability, while also acknowledging both its strengths in highlighting social cohesion and its limitations in addressing social change and conflict.
\item \textbf{Answer:} The United States government defines poverty through a poverty line established in the 1960 s, which quantifies the minimum income level necessary to secure basic necessities such as food, shelter, and clothing. Initially calculated based on the cost of a nutritionally adequate diet, this poverty threshold has been adjusted over the years to account for inflation, ensuring that it reflects contemporary economic conditions. However, these adjustments have often been criticized for not fully capturing the complexities of poverty, such as regional cost variations and non-monetary factors influencing well-being. Consequently, while the poverty line serves as a crucial indicator for social policy and resource allocation, it remains a limited measure that may overlook the multifaceted nature of poverty in modern society. The historical context of this definition highlights ongoing debates about the adequacy of government measures to address poverty in an evolving economic landscape.
\\ \textit{Note:} correct because it accurately describes the historical establishment of the poverty line in the 1960 s based on basic needs, acknowledges the adjustments for inflation over time, and recognizes the criticisms regarding the adequacy of these adjustments in reflecting current economic realities.
\end{enumerate}

\vfill
\noindent\textit{End of Answer Key}
\end{document}